\section{Introducción}
\label{sec:introduccion}

El presente informe consolida el proyecto de investigación desarrollado durante
el semestre bajo la metodología de Aprendizaje Basado en Proyectos, e integra en
un solo documento coherente los avances producidos en sus cuatro fases: la
delimitación del problema y la formulación de objetivos, la construcción del
marco teórico y la definición del enfoque metodológico, el diseño metodológico y
el levantamiento de datos, y finalmente la consolidación de resultados, su
discusión y las conclusiones del estudio.

El proyecto se inscribe en el ámbito de la inteligencia artificial aplicada a la
industria y de los productos tecnológicos orientados a la transformación
digital. Su objeto es el diseño, la codificación y la validación estadística de
un módulo de software analítico basado en el Filtro de Kalman Lineal Discreto,
concebido como el motor matemático de consistencia física de un gemelo digital
para el chancado primario minero. La relevancia disciplinar del trabajo no
reside en el algoritmo en sí, que es teoría establecida desde
\textcite{kalman1960}, sino en la posición que ese algoritmo ocupa dentro de una
arquitectura de decisión industrial: sostener que la inteligencia artificial
aplicada a la minería requiere obligatoriamente una capa previa de depuración
matemática de sus datos, porque ningún modelo predictivo puede ser más confiable
que la telemetría que lo alimenta.

La lógica de integración del informe sigue el recorrido investigativo real. Se
parte del fenómeno tecnológico observado —la degradación estocástica de las
señales SCADA en estaciones de chancado sometidas a vibración severa y polvo
abrasivo— y se lo delimita como un problema informático de observabilidad
parcial antes que como una falla de instrumentación. De esa delimitación se
deriva la pregunta de investigación, y de ella los objetivos, que ordenan el
trabajo en cuatro logros verificables: caracterizar el ruido, modelar la
dinámica, desarrollar el estimador y validarlo estadísticamente. La revisión de
literatura sitúa el estudio en la discusión académica vigente sobre gemelos
digitales industriales y estimación óptima, y de esa revisión emerge el vacío de
conocimiento que el proyecto aborda. El diseño metodológico traduce ese vacío en
un experimento computacional reproducible, cuyos resultados se presentan,
discuten y contrastan con los objetivos declarados.

Este documento se articula además con el artículo científico desarrollado en
paralelo, que reporta el mismo estudio en formato IEEE para una audiencia
especializada. El Anexo~\ref{anexo:correspondencia} explicita la correspondencia
sección a sección entre ambos productos, de modo que la trazabilidad entre el
informe y el manuscrito sea verificable y no meramente declarada.

Corresponde señalar de manera explícita que esta versión final incorpora
ajustes sustantivos derivados de la retroalimentación docente recibida durante
las fases previas. Se unificó la citación bajo normas APA 7, se completó la
formulación del plan del proyecto, se reestructuraron las consideraciones éticas
bajo el esquema riesgo--mitigación, se corrigió la coherencia dimensional de las
métricas de error, se verificó empíricamente el supuesto de normalidad de la
prueba de hipótesis y se depuró la bibliografía. El Anexo~\ref{anexo:ajustes}
documenta cada uno de estos ajustes, la observación que lo originó y la sección
del informe donde se materializa.

La estructura del documento es la siguiente. Las secciones \ref{sec:fenomeno}
a~\ref{sec:objetivos} presentan el fenómeno de estudio, la delimitación del
problema, la pregunta de investigación y los objetivos. La
sección~\ref{sec:estado-arte} desarrolla la revisión del estado del arte y
formula el vacío de conocimiento; la sección~\ref{sec:plan} presenta el plan del
proyecto con sus etapas, hitos y criterios de seguimiento; y la
sección~\ref{sec:marco} construye el marco teórico. Las
secciones~\ref{sec:enfoque} a~\ref{sec:analisis} exponen el diseño metodológico
completo. La sección~\ref{sec:resultados} reporta los resultados obtenidos, y
las secciones~\ref{sec:calidad} a~\ref{sec:coherencia} desarrollan los criterios
de calidad, los resguardos éticos y el análisis de coherencia interna del
diseño. Finalmente, las secciones~\ref{sec:discusion} y~\ref{sec:conclusiones}
interpretan los hallazgos y cierran el estudio.
